\subsection{Сферическое зеркало}

%%%%%%%%%%%%%%%%%%%%%%%%%%%%%%%%%%%%%%%%%%%%%%%%%%%%
%%%%%%%%%%%%%%%%%%%%%%%%%%%%%%%%%%%%%%%%%%%%%%%%%%%%
 \z{Найдите главное фокусное расстояние зеркала, если светящаяся точка и её действительное изображение лежат на главной оптической оси вогнутого зеркала на расстоянии 16 см и 100 см соответственно от главного фокуса.}
%
{$F=0.4\,\text{см}$}

%%%%%%%%%%%%%%%%%%%%%%%%%%%%%%%%%%%%%%%%%%%%%%%%%%%%
 \z{Радиус кривизны вогнутого зеркала $80$ см. На каком расстоянии от зеркала нужно поместить предмет, чтобы его действительное изображение было вдвое больше предмета?}
%
{$d=0.6\,\text{см}$}

%%%%%%%%%%%%%%%%%%%%%%%%%%%%%%%%%%%%%%%%%%%%%%%%%%%%
 \z{Предмет высотой $H = 5\,\text{см}$ находится на расстоянии $a = 12\,\text{см}$ от вогнутого зеркала с фокусным расстоянием $F = 10\,\text{см}$. На каком расстоянии от оптического центра зеркала находится изображение? Найти высоту изображения.}
%
{$b = 60\,\text{см}$, \quad $h = 25\,\text{см}$}

%%%%%%%%%%%%%%%%%%%%%%%%%%%%%%%%%%%%%%%%%%%%%%%%%%%%
 \z{Вогнутое сферическое зеркало даёт действительное изображение, которое в $\Gamma = 3$ раза больше предмета. Определить фокусное расстояние зеркала, если расстояние между предметом и его изображением $L = 20\,\text{см}$.}
%
{$F = 7.5\,\text{см}$}

%%%%%%%%%%%%%%%%%%%%%%%%%%%%%%%%%%%%%%%%%%%%%%%%%%%%
 \z{Каков радиус вогнутого сферического зеркала, находящегося на расстоянии $a = 2\,\text{м}$ от лица, если человек видит в нём своё изображение в $n = 1.5$ раза большим, чем в плоском зеркале, находящемся на том же расстоянии?}
%
{$R = 1.2\,\text{м}$}

%%%%%%%%%%%%%%%%%%%%%%%%%%%%%%%%%%%%%%%%%%%%%%%%%%%%
 \z[{\begin{wrapfigure}[8]{r}{0.4\textwidth}%
		\vspace{-20pt}%
		\begin{flushright}%
 			\includegraphics[width=0.4\textwidth]{4.2.6}%
		\end{flushright}%
	\end{wrapfigure}%
	}]%
	{Сходящиеся лучи падают на выпуклое зеркало так, что их продолжения пересекаются в точке, находящейся на расстоянии $l_1 = 0.4\,\text{м}$ за зеркалом (рис.). После отражения от зеркала лучи расходятся таким образом, что их продолжения пересекаются в точке, отстоящей от зеркала на расстоянии $l_2 = 1.6\,\text{м}$. Обе точки лежат на главной оптической оси зеркала. Определить фокусное расстояние зеркала.}
%
{$F=-32\,\text{см}$}

%%%%%%%%%%%%%%%%%%%%%%%%%%%%%%%%%%%%%%%%%%%%%%%%%%%%
 \z{Глаз наблюдателя расположен так, что муравей и его изображение в кривом (сферическом) зеркале для наблюдателя имеют одинаковые угловые размеры и полностью накладываются друг на друга. Наблюдатель отодвинулся от зеркала на расстояние $L$ вдоль линии, на которой находятся муравей и его изображение, и теперь видит, что угловой размер муравья составляет $75\%$ от углового размера его изображения. Затем наблюдатель отодвинулся в том же направлении ещё на $L$, и угловой размер изображения стал в $1.5$ раза больше углового размера муравья. Во сколько раз изображение муравья больше его самого? Каков радиус кривизны зеркала?}
%
{В два раза, \quad $R = \frac{4}{3}L$}


